\section{跨角度、频率和偏振的统一验证}

\subsection{高阶微分测试相同的形式体系是否能处理更锐利的算子曲率}

如果提出的响应空间表述真正比任务特定配方更基础，那么当目标曲率改变而底层逆向设计机制不变时，它应保持连贯。四阶微分结果提供了该测试。在项目档案中，四阶候选使用与二阶基准采用的相同通用逻辑进行评分，但目标曲线对应于对横向波矢量更陡的依赖性，并显式跟踪中心抑制、形状匹配、边缘行为和有效带宽。多个波长上可行的高分候选的存在表明，当前设计策略不限于一个预选传递函数轮廓。

在科学上，这很重要，因为高阶微分不是二阶情况的平凡扰动。更锐利的目标曲率对超表面响应和排序标准都提出更强要求。名义上有希望的设计更容易匹配中心而在边缘失败，或表现出正确的定性对称性而错过有用的算子保真度所需的定量角度增长。因此，四阶结果探测被学习的响应流形是否足够广泛以包括多个算子阶数，而非一个偏好的响应族。仓库输出建议的答案是肯定的：当算子阶数改变时，相同的候选生成和物理重排序程序保持有意义。

\subsection{低通和高通目标测试光谱-角度滤波是否能被纳入相同的响应语言}

角度低通和高通任务探测统一表述的不同轴。它们不改变算子类目标的多项式阶数，而是重塑角度包络本身。仓库用搜索合理有效角度宽度然后评估中心行为、形状一致性和带宽支持的分数函数实现两个目标族。这很重要，因为它将可能作为新应用呈现的内容转换为更简单的概念主张：角度滤波是响应空间中的另一种目标几何。

结果支持这种解读。在低通摘要中，高分候选对应于保留强中心区域同时拒绝外角度带的响应；在高通摘要中，权重有效反转，有利于中心抑制和强外角度响应。两种情况下使用相同的响应输入、候选生成和RCWA重排序程序。不引入任务特定的逆向架构，几何表示保持不变。这些观察强化了本文的中心叙述。光谱或角度滤波不需要被框定为添加到基线微分器配方的"额外任务"。它已经是响应空间表述的原生内容，因为它由相同角度--波长坐标上的不同目标包络定义。

\subsection{偏振无关和偏振复用设计测试偏振是否作为一等设计轴进入}

偏振相关任务针对统一表述的最清晰可能挑战。在许多超表面工作流程中，偏振作为可选扩展或单独设计分支被后期引入。这里它从一开始就被构建进响应张量。仓库包括偏振无关目标（奖励匹配通道行为）和偏振复用目标（奖励一个通道中的活动与另一个通道中的抑制并存）。分数摘要显式报告联合分数、通道特定形状分数、无源通道泄漏以及大正负角度处的通道值。这正是如果偏振被作为响应对象的平等部分而非附加约束处理所预期的记账类型。

偏振无关情况询问相同的设计策略是否能在目标角度窗口上保持两个通道对齐。偏振复用情况询问更难的问题：相同的逆向表述能否在一个通道保持功能活动而另一个被驱动向静默的结构中产生？从验证摘要中出现的答案是物理上合理的方式细化的。强候选确实存在，但主动通道形状匹配与无源通道抑制之间的差距揭示了真正的工程权衡而非平凡的设计条件。这种权衡正是物理有意义的表述应该揭示的。将其隐藏在单一平均分数后面的响应空间模型的信息量会更少，而非更多。

综合起来，高阶、滤波和偏振目标结果支持相同的结论。贡献不在于一个代码路径可以指向几个不相关的算子任务。贡献在于这些任务可以在一个角度--频率--偏振设计语言内表示、采样和验证。因此，统一在算法之前是概念性的：响应流形是正确的学习对象，跨目标结果显示该对象足够丰富以容纳多个实际相关的傅里叶算子行为。
